\documentclass[12pt]{article}

\usepackage[utf8]{inputenc}
\usepackage[T1]{fontenc}
\usepackage[margin=1in]{geometry}
\usepackage[pdftex]{hyperref}
\usepackage{amsmath,amsthm,amssymb,graphicx,mathtools,tikz,hyperref,enumerate}
\usepackage{mdframed,cleveref,cancel,stackengine,pgfplots,pgf,mathrsfs,thmtools}
\usepackage{xfrac,stmaryrd,commath,needspace,multirow}
\usepackage[catalan]{babel}

\newmdenv[leftline=false,topline=false]{topright}
\let\proof\relax
\usetikzlibrary{positioning,arrows, calc, babel}
\usetikzlibrary{external}
\tikzexternalize[prefix=figures/]
\pgfplotsset{compat=1.11}

\def\fbx#1{\vbox{\hbox{\hbox{#1}\setbox0\lastbox\copy0\kern\fboxsep\vrule width\fboxrule depth\dimexpr \fboxsep+\dp0\relax}%
            \hrule height\fboxrule}}

\newcommand{\bimplies}{\fbx{$\implies$}}
\newcommand{\bimpliedby}{\fbx{$\impliedby$}}

\newcommand{\n}{\mathbb{N}}
\newcommand{\z}{\mathbb{Z}}
\newcommand{\q}{\mathbb{Q}}
\newcommand{\cx}{\mathbb{C}}
\newcommand{\real}{\mathbb{R}}
\newcommand{\E}{\mathbb{E}}
\newcommand{\F}{\mathbb{F}}
\newcommand{\A}{\mathbb{A}}
\newcommand{\R}{\mathcal{R}}
\newcommand{\C}{\mathcal{C}}
\newcommand{\Pa}{\mathcal{P}}
\newcommand{\Es}{\mathcal{E}}
\newcommand{\V}{\mathcal{V}}
\newcommand{\bb}[1]{\mathbb{#1}}
\newcommand{\pdv}[3][]{\frac{\partial^{#1} #2}{\partial #3^{#1}}}
\newcommand{\dv}[3][]{\frac{\dif^{#1} #2}{\dif #3^{#1}}}
\let\k\relax
\newcommand{\k}{\Bbbk}
\newcommand{\ita}[1]{\textit{#1}}
\newcommand\inv[1]{#1^{-1}}
\newcommand\setb[1]{\left\{#1\right\}}
\newcommand{\vbrack}[1]{\langle #1\rangle}
\newcommand{\determinant}[1]{\begin{vmatrix}#1\end{vmatrix}}
\newcommand{\Po}{\mathbb{P}}
\newcommand{\lp}{\left(}
\newcommand{\rp}{\right)}
\DeclareMathOperator{\fr}{Fr}
\DeclareMathOperator{\Id}{Id}
\DeclareMathOperator{\ext}{Ext}
\DeclareMathOperator{\inte}{Int}
\DeclareMathOperator{\rie}{Rie}
\DeclareMathOperator{\rg}{rg}
\DeclareMathOperator{\gr}{gr}
\DeclareMathOperator{\nuc}{Nuc}
\DeclareMathOperator{\car}{car}
\DeclareMathOperator{\im}{Im}
\DeclareMathOperator{\tr}{tr}
\DeclareMathOperator{\vol}{vol}
\DeclareMathOperator{\grad}{grad}
\DeclareMathOperator{\rot}{rot}
\DeclareMathOperator{\diver}{div}
\DeclareMathOperator{\sinc}{sinc}
\DeclareMathOperator{\graf}{graf}
\DeclareMathOperator{\tq}{\;t.q.\;}
\DeclareMathOperator{\disc}{disc}
\let\emptyset\varnothing
\renewcommand{\thesubsubsection}{\Alph{subsubsection}}
\setcounter{secnumdepth}{4}

\hypersetup{
    colorlinks,
    linkcolor=blue
}
\def\upint{\mathchoice%
    {\mkern13mu\overline{\vphantom{\intop}\mkern7mu}\mkern-20mu}%
    {\mkern7mu\overline{\vphantom{\intop}\mkern7mu}\mkern-14mu}%
    {\mkern7mu\overline{\vphantom{\intop}\mkern7mu}\mkern-14mu}%
    {\mkern7mu\overline{\vphantom{\intop}\mkern7mu}\mkern-14mu}%
  \int}
\def\lowint{\mkern3mu\underline{\vphantom{\intop}\mkern7mu}\mkern-10mu\int}

\renewcommand*\contentsname{Contenidos}

\newtheoremstyle{break}% name
{}%         Space above, empty = `usual value'
{}%         Space below
{}% Body font
{}%         Indent amount (empty = no indent, \parindent = para indent)
{\bfseries}% Thm head font
{}%        Punctuation after thm head
{\newline}% Space after thm head: \newline = linebreak
{#1 #2 \normalfont {\it #3}}%         Thm head spec

\newtheoremstyle{breakthm}% name
{}%         Space above, empty = `usual value'
{}%         Space below
{}% Body font
{}%         Indent amount (empty = no indent, \parindent = para indent)
{\bfseries}% Thm head font
{}%        Punctuation after thm head
{\newline}% Space after thm head: \newline = linebreak
{#1 \normalfont #3 (#2)\addcontentsline{toc}{subsection}{#1 #3}}%         Thm head spec
\newtheoremstyle{normal}% name
{}%         Space above, empty = `usual value'
{}%         Space below
{}% Body font
{}%         Indent amount (empty = no indent, \parindent = para indent)
{\bfseries}% Thm head font
{}%        Punctuation after thm head
{5pt plus 1pt minus 1pt}% Space after thm head: \newline = linebreak
{#1 #2 \normalfont #3}%         Thm head spec

\theoremstyle{normal}
\declaretheorem[style=normal,name=Lema,numberwithin=section]{lema}
\declaretheorem[style=normal,name=Lema,numbered=no]{lema*}
\declaretheorem[style=normal,name=Observació,sibling=lema]{obs}
\declaretheorem[style=normal,name=Observació,numbered=no]{obs*}

\declaretheorem[style=break,name=Proposició,sibling=lema]{prop}
\declaretheorem[style=break,name=Proposició,numbered=no]{prop*}
\declaretheorem[style=break,name=Demostració,qed=$\square$,numbered=no]{proof}
\declaretheorem[style=break,name=Definició,sibling=lema]{defi}
\declaretheorem[style=break,name=Definició,numbered=no]{defi*}
\declaretheorem[style=break,name=Corol·lari,sibling=lema]{col}
\declaretheorem[style=break,name=Corol·lari,numbered=no]{col*}
\declaretheorem[style=break,name=Exercici,sibling=lema]{ej}
\declaretheorem[style=break,name=Exercici,numbered=no]{ej*}
\declaretheorem[style=break,name=Exemple,sibling=lema]{example}
\declaretheorem[style=break,name=Exemple,numbered=no]{example*}

\let\thm\relax
\declaretheorem[style=breakthm,name=Teorema,sibling=lema]{teo}
\declaretheorem[style=break,name=Teorema,sibling=lema]{teo*}
\declaretheorem[style=breakthm,name=M\`etode,sibling=lema]{met}

\begin{document}

\newcommand{\plogo}{\fbox{$\mathcal{PL}$}} % Generic dummy publisher logo

\begin{titlepage} % Suppresses headers and footers on the title page

	\centering % Centre everything on the title page

	\scshape % Use small caps for all text on the title page

	\vspace*{\baselineskip} % White space at the top of the page

	%------------------------------------------------
	%	Title
	%------------------------------------------------

	\rule{\textwidth}{1.6pt}\vspace*{-\baselineskip}\vspace*{2pt} % Thick horizontal rule
	\rule{\textwidth}{0.4pt} % Thin horizontal rule

	\vspace{0.75\baselineskip} % Whitespace above the title

	{\LARGE CÁLCULO INTEGRAL\\} % Title

	\vspace{0.75\baselineskip} % Whitespace below the title

	\rule{\textwidth}{0.4pt}\vspace*{-\baselineskip}\vspace{3.2pt} % Thin horizontal rule
	\rule{\textwidth}{1.6pt} % Thick horizontal rule

	\vspace{2\baselineskip} % Whitespace after the title block

	%------------------------------------------------
	%	Subtitle
	%------------------------------------------------

	%A Number of Fascinating and Life-changing Templates Presented in a Clear and Concise Way % Subtitle or further description

	\vspace*{3\baselineskip} % Whitespace under the subtitle

	%------------------------------------------------
	%	Editor(s)
	%------------------------------------------------



	\vspace{1.5\baselineskip} % Whitespace before the editors

	{\scshape\Large Òscar Benedito \\ Jordi Castellví \\ Ernesto Lanchares \\ Miquel Ortega \\} % Editor list

	\vspace{3\baselineskip} % Whitespace below the editor list

	\textit{Universitat Politècnica de Catalunya \\ Barcelona} % Editor affiliation

	\vfill % Whitespace between editor names and publisher logo

	%------------------------------------------------
	%	Publisher
	%------------------------------------------------

	%\plogo % Publisher logo
	\includegraphics[width=5cm]{../logo.pdf}

	\vspace{0.3\baselineskip} % Whitespace under the publisher logo
        
        % \date{\ifcase \month \or Enero\or Febrero\or Marzo\or Abril\or Mayo \or Junio\or Julio\or Agosto\or Setiembre\or Octubre\or Noviembre\or Diciembre\fi \:  \number \year} \let\Date\@date

	Diciembre 2017 % Publication year

	%{\large ApuntsFME} % Publisher

\end{titlepage}

\section{Programaci\'o no lineal sense restriccions}

\begin{teo*}[Condicions necess\`aries d'optimalitat]
    Sigui $f \colon \real^n \to \real$ un problema d'optimitzaci\'o no lineal $\min_{x\in\real^n} f\left( x \right)$. Si $x^*$ \'es un m\'inim local de $f$ i $f \in \mathcal{C}^2$ en un entorn de $x^*$, llavors:
    \begin{enumerate}[i)]
        \item $\nabla f\left( x^* \right) = 0$. (Condici\'o de 1r ordre)
        \item $\nabla^2 f\left( x^* \right) \geq 0$. (Condici\'o de 2n ordre)
    \end{enumerate}
\end{teo*}
\begin{teo*}[Condicions suficients d'optimalitat]
    Si $f \in \mathcal{C}^2$ en un entorn obert de $x^*$, $\nabla f\left( x^* \right) = 0$ i $\nabla^2 f\left( x^* \right)$ \'es definida positiva, llavors $x^*$ \'es m\'inim local estricte de $f$.
\end{teo*}
\begin{teo*}[Condicions d'optimalitat en problemes convexos]
    Si $f$ \'es convexa i diferenciable, llavors $\nabla f\left( x^* \right) = 0 \iff x^*$ \'es m\'inim global de $f$.
\end{teo*}
\begin{met}[del Gradient]
    $x^{k+1} = x^k + \alpha^kd^k$, \quad $d^k = -\nabla f\left( x^k \right), \quad \alpha$ ha de complir les condicions d'Armijo-Wolfe.
\end{met}
\begin{met}[de Newton]
    $x^{k+1} = x^k + \alpha^kd^k$, \quad $d^k = -\left( \nabla^2 f\left( x \right) \right)^{-1} \nabla f\left( x^* \right)$, \quad$\alpha^k = 1$.
\end{met}
\begin{prop}[Condicions d'Armijo-Wolfe]
    \begin{enumerate}[i)]
        \item[]
        \item Condici\'o de descens suficient (AW-1): \\
            \indent $g\left( \alpha \right) \leq g\left( 0 \right) + \alpha c_1g'\left( 0 \right), \quad c_1 \in (0,1).$ \\
            \indent $f\left( x^k + \alpha d^k \right) \leq f\left( x^k \right) + \alpha c_1 \nabla f\left( x^k \right)^t d^k, \quad c_1 \in (0,1).$
        \item Condici\'o de corbatura (AW-2): \\
            \indent $g'\left( \alpha \right) \geq c_2 g'\left( 0 \right), \quad 0 < c_1 < c_2 < 1.$ \\
        \indent $\nabla f\left( x^k + \alpha d^k \right)^td^k \geq c_2 \nabla f\left( x^k \right)d^k, \quad 0 < c_1 < c_2 < 1$.
    \end{enumerate}
\end{prop}

\section{Programaci\'o no lineal amb restriccions}

Donat el problema:
\begin{equation*}
    \lp P \rp \begin{dcases*}
        \min f \left( x \right) \\
        \text{s.a.: } \\
        \quad \left( 1 \right) \quad h \left( x \right) = 0 \\
        \quad \left( 2 \right) \quad g \left( x \right) \leq 0 \\
    \end{dcases*}
\end{equation*}
amb $\mathcal{L} \left( x, \lambda, \mu \right) = f\left( x \right) + \lambda^t h\left( x \right) + \mu^t g\left( x \right)$.

\begin{prop}[Condicions necess\`aries]
    Sigui $x^*$ \`optim local. Si \'es punt regular, llavors $\exists \lambda^* \in \real^m$ i $\mu^* \in \real^p$ tals que:
    \begin{enumerate}[i)]
        \item \label{nec:1} $h\left( x^* \right) = 0,\, g\left( x^* \right) \leq 0$.
        \item \label{nec:2} $\nabla_x \mathcal{L}\left( x^*, \lambda^*, \mu^* \right) = \nabla f\left( x^* \right) + \nabla h\left( x^* \right) \lambda^* \nabla g\left( x^* \right)\mu^* = 0$.
        \item \label{nec:3} $\mu^* \geq 0$ i $\left( \mu^* \right)^t g\left( x^* \right) = 0$ (si $g_i\left( x^* \right)$ \'es inactiva, llavors $\mu_j^* = 0$).
        \item \label{nec:4} $d^t \nabla_{x_ix_j}^2 \mathcal{L} \left( x^*, \lambda^*, \mu^* \right) d \geq 0, \forall d\in M =
            \begin{cases}
                \left( \nabla h_i\left( x^* \right) \right)^t d = 0 & i \in \{1,\dots,m\} \\
                \left( \nabla g_i\left( x^* \right) \right)^t d = 0 & j \in \mathcal{A}\left( x^* \right)
            \end{cases}$
    \end{enumerate}
    Els punts \ref{nec:1}, \ref{nec:2} i \ref{nec:3} s\'on les condicions de 1r ordre del KKT i el punt \ref{nec:4} \'es la condici\'o de 2n ordre.
\end{prop}
\begin{prop}[Condicions suficients]
    Sgui $x^*$, \'es \`optim local si satisf\`a:
    \begin{enumerate}[i)]
        \item \label{suf:1} $h\left( x^* \right) = 0,\, g\left( x^* \right) \leq 0$.
        \item \label{suf:2} $\nabla_x \mathcal{L} \left( x^*, \lambda^*, \mu^* \right) = \nabla f\left( x^* \right) + \nabla h\left( x^* \right)\lambda^* + \nabla g\left( x^* \right)\mu^* = 0$.
        \item \label{suf:3} $\mu^* \geq 0$ i $\left( \mu^* \right)^t g\left( x^* \right) = 0$ ($g_i\left( x^* \right) < 0 \implies \mu_j^* = 0$).
        \item \label{suf:4} $d^t \nabla_{x_ix_j}^2 \mathcal{L}\left( x^*, \lambda^*, \mu^* \right) d > 0, \quad d \in M$.
    \end{enumerate}
    Els punts \ref{suf:1}, \ref{suf:2} i \ref{suf:3} s\'on les condicions de 1r ordre del KKT i el punt \ref{suf:4} \'es la condici\'o de 2n ordre.
\end{prop}


\end{document}
